\section{AA制和抢单}
AA制源于英文“go dutch”，意思是个人平均分摊所需费用，通常用于饮食聚会及旅游等共同消费的场合，人们在消费后分摊消费总额，互不相欠。普遍的说法是，AA制起源于荷兰，在16-17世纪荷兰和威尼斯的海上贸易十分发达，意大利人与荷兰人在各个地方来回奔波，流动性很大，为了避免各自吃亏，在他们中间就产生了聚时交流信息，散时各付资费的传统。
而抢单，顾名思义，意为抢着买单，与AA制不同，通常是在吃晚饭结账时，一桌人抢着买单。更有甚者，为抢占先机，锅里饭菜还没吃完，就找个上个厕所的借口，偷偷到柜台处结账买单。
在西方社会，在亲朋好友之间的聚会中，人们实行AA制是非常普遍的。但是在中国，AA制一般只被青年人，特别是大学生和白领所接受，传统的中国中老年人相较于AA制，更习惯于在聚会时替人买单。
究其原因，一方面西方人的价值观更偏向于利益至上。上至国家，下至个人，利益至上是西方社会普遍认同的价值观。因此，西方人上餐厅吃饭，就算两人感情再好，吃完了各付个的，这非常的自然，没有人会认为这不妥。只有这样，才能体现对别人和对自己私有财产的尊重，才能保全各自在金钱上的利益。而与之相反的是，尽管中国人普遍接受了派对、各种西方节日等丰富的西方生活方式，但AA制却始终未能流行开来，原因之一就是AA制与中国人看重人情、面子的文化传统相违背。看重情义、面子是中国人的性格特点。这种特点体现在对父母，对兄弟，师长，朋友等方方面面中。对中国人来说，很难接受与朋友出去吃饭，还要各付各的。这样子做会显得自己吝啬，也会被周遭人侧目相看。另外，中国人好客，爱请客。如果人们热热闹闹地围着一张大桌子，欢谈畅饮，最后结账时却分摊消费，这对中国人来说是丢面子的事，而面子，正是中国人最看重的东西。所以，许多国人宁愿打肿脸充胖子，也要抢着付钱，积极程度可与打架相当。
\begin{figure}
    \centering
    \includegraphics[width=0.6\linewidth]{sammelauftraege}
    \caption{在中国抢单}
\end{figure}
\par
另一方面，在中国，吃饭有着独特的文化内涵。中国人喜欢把事情搬到饭桌上解决。求人办事时会请人吃饭，答谢他人时会请人吃饭，送行时也会请人吃饭，如此这般。关于中国人请客吃饭时抢单的原因，一方面是因为纯粹的感情，另一方面也是因为人们为了办事而不得不进行应酬。请客在中国实际上是一种长期的投资，除了投入金钱外，也投入了长期的人情。这与西方公平保证双方的利益的AA制有着本质的区别。一般来说，被请客之人往往会有一种吃人嘴软的感觉，觉得欠了别人人情，以后被求帮忙也就不好意思拒绝了。人情是中国人实现利益最重要的手段之一，在中国，即使人人都明白这个道理，也却一次次掉入这个陷阱中，让请客之人百试不爽。而相对的，西方人不喜欢欠他人人情，事情必须一件件的算清楚，他们不觉得AA制让人尴尬，若你有事求他，而他不能从中得到好处，他可以很干脆地回绝你，因为在他们眼中，只有利益是最重要的，人情，面子，都是虚的。
而抢单文化在中国也并非是一伙人吃饭每次到了结账关头次次抢单，通常是这一次我请大家吃一顿，下一次你回请我吃一顿。这样子既有了面子，也有了利益的保障。在实际生活中，可能存在着有些朋友，他们家庭条件不好，但是感情到位了，大家就一直不让他买单，这也是普遍的事，真正的朋友之间，谁也不会在乎这点饭钱，抢单文化在这里也不再是对于人情的争抢，而是朋友之间的关怀与暖意。